% Generated draft: requires mathematical and source review.
\documentclass[11pt]{article}
\usepackage[margin=25mm]{geometry}
\usepackage{fontspec}
\setmainfont{Latin Modern Roman}
\usepackage{amsmath,amssymb,mathtools,unicode-math}
\setmathfont{Latin Modern Math}
\usepackage{enumitem,needspace,xurl,hyperref}
\setlength{\parindent}{0pt}
\setlength{\parskip}{5pt}
\setlength{\emergencystretch}{3em}
\newcommand{\N}{\mathbb{N}}
\newcommand{\Z}{\mathbb{Z}}
\newcommand{\Q}{\mathbb{Q}}
\newcommand{\R}{\mathbb{R}}
\newcommand{\C}{\mathbb{C}}
\newcommand{\F}{\mathbb{F}}
\newcommand{\A}{\mathbb{A}}
\newcommand{\PP}{\mathbb{P}}
\newcommand{\E}{\mathbb{E}}
\newcommand{\eps}{\varepsilon}
\newcommand{\dd}{\,\mathrm{d}}
\begin{document}
\section*{UnsolvedMath ID: 30001126}
\textbf{Source problem code:} catalogue entry
\subsection{Definitions and mathematical statement}
A standard framework for interpreting the question is a phase space \(X\) carrying a Hamiltonian evolution \(\Phi_t\), together with a Lie group \(G\) acting by symmetries, so that
\[
\Phi_t(g\cdot u)=g\cdot\Phi_t(u)
\]
whenever the evolution is defined. This framework explains the terminology; the supplied record does not specify \(X\), its Hamiltonian or Poisson structure, the symmetry group, or the evolution equation.

In this framework, a solution through \(u_0\in X\) is a \emph{relative equilibrium} if there exists \(\xi\) in the Lie algebra of \(G\) such that
\[
\Phi_t(u_0)=\exp(t\xi)\cdot u_0
\]
for every time in its interval of existence. A solution \(u(t)\) is \emph{relative periodic} if there exist \(T>0\) and \(g\in G\) such that
\[
u(t+T)=g\cdot u(t)
\]
whenever both sides are defined. Relative equilibria can also satisfy this relative-periodicity condition; requiring a relative periodic solution that is not a relative equilibrium is an additional distinction whose precise role must be recovered from the underlying transition results.

\textbf{Source-level problem.} Can results on the Hamiltonian meandering transition---from relative equilibria to relative periodic orbits---be extended to infinite-dimensional Hamiltonian systems, such as PDE models of vortex dynamics?

Here ``Hamiltonian meandering transition'' retains the source terminology. The supplied record does not identify the finite-dimensional results to be extended, define the transition precisely, specify a parameterized family, or provide hypotheses and a proposed conclusion for an infinite-dimensional theorem. PDE models of vortex dynamics are given as examples, without a specified equation, spatial domain, boundary conditions, or solution class. Consequently, the record determines a research question but does not determine a unique fully quantified extension theorem. These qualifications remain unresolved in this statement-only draft.
\subsection{Short English statement}
The problem asks whether Hamiltonian results describing transitions from relative equilibria to relative periodic motion extend to infinite-dimensional systems, including PDE models of vortex dynamics. The specific results, equations, and hypotheses are unspecified, so a precise theorem requires further source clarification. The supplied status is \emph{partially solved}; it is not independently verified here.
\subsection{Sources}
\begin{itemize}
\item Supplied local UnsolvedMath registry record, permanent ID \(30001126\), source problem code \texttt{catalogue entry}, title ``Hamiltonian Meandering Transitions in Infinite Dimensions,'' in the collection ``Oberwolfach Reports Open Problems.'' No individual proposer is identified in the supplied record; the recorded proposed year is 2008.
\item Supplied source citation and literature reference ``original OWR report'': \emph{Dynamics of Patterns (2009)}, Oberwolfach Reports, pp.~3201--3271, DOI \texttt{10.4171/owr/2008/57}. \url{https://doi.org/10.4171/owr/2008/57}. The background describes Oberwolfach Reports as reports of mathematical workshops and their open-problem discussions. The citation's 2009 label and the record's 2008 proposed year are retained without resolving their relationship.
\item Supplied literature reference ``KAM for vortex patches'': \url{https://doi.org/10.1134/S1560354724540013}. No authors, theorem statement, or explanation of its precise implications for this transition is supplied.
\item Literature assessment dated exactly \texttt{2026-08-21}: the supplied curation labels the problem \emph{partially solved} and reports substantial KAM and quasi-periodic-solution theory for infinite-dimensional Hamiltonian PDEs, including vortex-patch results, while describing a general extension to PDE models of vortex dynamics as open. This is supplied dated curation derived from the OWR record, not an independently verified literature assessment; its time-sensitive status and cited claims require independent verification.
\item The registry labels the statement \texttt{exact} and reports that its mathematical content was checked against the official Oberwolfach report, with notation, context, and question form normalized without changing the source-level claim. This is a supplied verification assertion; no independent source consultation is claimed here.
\end{itemize}
\end{document}
